% Target length in the official two-column template: approximately 1.5 pages.
\section{Methodology}
\label{sec:method}

\subsection{Routing Objective}

Let $x$ be an input tuple, $q$ a semantic operator, and $M_N,M_M$ the Nano and
Mini answer models. Their predictions are $\hat{z}_N(x,q)$ and
$\hat{z}_M(x,q)$. A pre-router $r_\theta(x,q)\in\{N,M\}$ executes before either
answer model and minimizes expected execution cost subject to a quality target:
\begin{equation}
  \min_\theta\; \mathbb{E}[c_R(x,q)+c_{r_\theta(x,q)}(x,q)]
  \quad\text{s.t.}\quad
  A(r_\theta) \ge A(M_M)-\epsilon,
  \label{eq:constrained}
\end{equation}
where $c_R$ is routing overhead, $A$ is end-to-end answer accuracy, and
$\epsilon$ is an allowed loss. We do not solve Eq.~\ref{eq:constrained}
directly. Training optimizes a scalar surrogate, and a validation-only
candidate/threshold sweep constructs an empirical cost--accuracy frontier. For
an operating point claimed under Eq.~\ref{eq:constrained}, we select the
cheapest validation candidate satisfying the accuracy floor. Historical rows
retain their original calibration rules and are not claimed as direct
solutions of Eq.~\ref{eq:constrained}.

\subsection{Cost-Derived Score Function}

For labeled offline data, paired execution defines outcome
$b(x)=(\mathbb{1}[\hat z_N=z],\mathbb{1}[\hat z_M=z])$, yielding both correct
(\bc), Mini only (\mo), Nano only (\no), or both wrong (\bw). Each cached answer
also records its input and output token counts. Given the price snapshot
$\pi_r^{\mathrm{in}},\pi_r^{\mathrm{out}}$ per million tokens, realized answer
cost is
\begin{equation}
 c_r(x,q)=\frac{\pi_r^{\mathrm{in}}n_{r,x}^{\mathrm{in}}+
                    \pi_r^{\mathrm{out}}n_{r,x}^{\mathrm{out}}}{10^6}.
 \label{eq:realized-cost}
\end{equation}
We estimate each route's expected cost on the training split,
\begin{equation}
 \begin{aligned}
 \bar c_r&=\mathbb{E}_{x\sim\mathcal D_{\mathrm{train}}}[c_r(x,q)]\\
 &=\frac{1}{|\mathcal D_{\mathrm{train}}|}
   \sum_{x\in\mathcal D_{\mathrm{train}}}c_r(x,q)
 \end{aligned}
 \label{eq:expected-cost}
\end{equation}
and use its inverse as the common cost scale:
\begin{equation}
 S_{\mathrm{cost}}(x,r)=
 \frac{\mathbb{1}[\hat z_r(x,q)=z]}{\bar c_r}.
 \label{eq:cost-score}
\end{equation}
A wrong answer receives zero; a correct answer from a lower-expected-cost route
receives a larger score. Multiplying all prices by a common currency conversion
changes score scale but not route ordering.

\begin{table}[t]
  \centering
  \small
  \caption{Cost-derived route score using expected route costs.}
  \label{tab:score}
  \begin{tabular}{lccc}
    \toprule
    Outcome & $S(x,N)$ & $S(x,M)$ & Preferred route\\
    \midrule
    \bc & $1/\bar c_N$ & $1/\bar c_M$ & lower-cost correct model\\
    \mo & $0$ & $1/\bar c_M$ & Mini\\
    \no & $1/\bar c_N$ & $0$ & Nano\\
    \bw & $0$ & $0$ & indifferent; Nano tie-break\\
    \bottomrule
  \end{tabular}
\end{table}
Table~\ref{tab:score} makes the score semantics explicit without a hand-tuned
reward table. A critical under-escalation on \mo forfeits all recoverable
score, while sending \bc to Mini remains correct but earns less whenever
$\bar c_N<\bar c_M$. Neither route can recover \bw; the deployment hard-label rule uses
Nano and reports \bw separately. The \emph{conservative accuracy oracle} in
Tables~\ref{tab:historical-main} and~\ref{tab:fresh} intentionally overrides
this tie and sends \bw to Mini. That convention supplies an accuracy ceiling
with conservative escalation; it is neither the cost-optimal oracle nor the
recommended deployment behavior, and the \bw choice cannot change accuracy.
Invalid discrete output receives zero score and falls back to Mini
operationally. The recorded QA runs use the pricing snapshot embedded in their
artifacts: GPT-4.1 Nano at \$0.10/\$0.40 and Mini at \$0.40/\$1.60 per million
input/output tokens.

\subsection{Disagreement-Guided Critical-Case Acquisition}
\label{sec:active-acquisition}

Let $\mathcal U=\{(x,q)\}_{i=1}^{N}$ be a candidate pool with no revealed gold
answers. We first execute both answer models on every candidate and cache
$\hat z_N(x,q)$ and $\hat z_M(x,q)$. After applying the task's fixed answer
normalizer, these observable predictions define
\begin{equation}
 \begin{aligned}
 \mathcal A &= \{x:\hat z_N(x,q)=\hat z_M(x,q)\},\\
 \mathcal D &= \{x:\hat z_N(x,q)\ne\hat z_M(x,q)\},
 \end{aligned}
 \label{eq:agreement-strata}
\end{equation}
where $\mathcal A$ and $\mathcal D$ are the agreement and disagreement strata.
For exact-label tasks equality is literal; a free-form operator must declare
its equivalence evaluator before acquisition.

Given a human-label budget $B\ll N$, we choose
$\mathcal S_B=\mathcal S_D\cup\mathcal S_A$, with
$\mathcal S_D\subseteq\mathcal D$, $\mathcal S_A\subseteq\mathcal A$, and
$|\mathcal S_D|+|\mathcal S_A|=B$. The policy prioritizes disagreements and
uses a smaller random agreement sample for coverage. Only after
$\mathcal S_B$ is fixed do annotators reveal gold $z$ for those tuples.
Consequently, neither selection nor its tie-breaking can inspect correctness or
the latent \bc/\mo/\no/\bw bucket. Gold and the cached pair then determine the
route target and score in Eq.~\ref{eq:cost-score}.

Disagreement is useful because every \mo and \no case must lie in
$\mathcal D$: one candidate model matches gold and the other does not. These
are the examples that teach when escalation helps and when Nano is preferable.
The disagreement stratum is not itself a correctness label, because \bw can
also disagree. Conversely, agreement examples include ordinary \bc cases that
teach safe Nano routing, as well as shared \bw errors. Retaining agreement
coverage therefore controls over-escalation and exposes systematic shared
failures.

Let $h(x)$ denote the cost of obtaining a human gold answer. Full annotation
and active annotation cost
\begin{align}
 C_{\mathrm{gold}}^{\mathrm{full}} &= \sum_{x\in\mathcal U}h(x),\\
 C_{\mathrm{gold}}^{\mathrm{active}} &= \sum_{x\in\mathcal S_B}h(x).
 \label{eq:annotation-cost}
\end{align}
With constant annotation cost, the gold-label saving is exactly $1-B/N$.
Both active and random selection in our matched comparison pay the same paired
screening cost $\sum_{x\in\mathcal U}(c_N(x)+c_M(x))$. The method therefore
claims human-annotation efficiency, not a reduction in Nano or Mini screening
calls. Its value is that the limited gold budget contains more routing-critical
supervision.

\subsection{Primary Lightweight Soft-Prompt Router}

The primary router uses a frozen encoder rather than a generative LLM. Let
$E(x,q)$ be DistilBERT embeddings and
$V\in\mathbb{R}^{p\times d}$ a trainable soft prompt. The input
$[V;E(x,q)]$ passes through the frozen encoder, and a two-class linear head
produces $p_\theta(r\mid x)$. Only $V$ and the head are updated. For
$d=768$ and $p\in\{2,4,8,16\}$, this is 3,074--13,826 trainable parameters.
The task-specific state is therefore tiny, and the fixed encoder is also
substantially smaller than the separate 2B discrete router. Inference is one
encoder pass with no autoregressive route-token generation.

Active acquisition chooses which paired tuples train the router. The cost score
defines the hard target and the weight for every selected example:
\begin{equation}
 y_i=\arg\max_{r\in\{N,M\}}S_{\mathrm{cost}}(x_i,r),
 \qquad
 w_i=\frac{1}{\bar c_{y_i}},
 \label{eq:cost-target-weight}
\end{equation}
with Nano used when the route scores tie. The soft prompt is trained only with
inverse-expected-cost weighted cross-entropy,
\begin{equation}
 \mathcal{L}_{\mathrm{WCE}}=
 \frac{\sum_i w_i\,
 \ell_{\mathrm{CE}}\!\left(p_\theta(\cdot\mid x_i),y_i\right)}
 {\sum_i w_i}.
 \label{eq:weighted-ce}
\end{equation}
There are no hand-tuned outcome weights and no auxiliary differentiable loss.
Because inverse-cost weighting alone does not upweight rare \mo cases, the
disagreement-guided acquisition stage is responsible for placing enough critical cases
in the training set. This separates two roles cleanly: acquisition controls
rare-case coverage, while Eq.~\ref{eq:weighted-ce} applies one auditable cost
weighting rule to all examples.

Validation enumerates thresholds over $p_\theta(M\mid x)$, selects maximum
answer accuracy with fewer Mini calls as the tie-break, and reports critical
misroutes and realized cost. For Eq.~\ref{eq:constrained}, we instead select the
cheapest validation threshold satisfying the accuracy floor.

\subsection{Independent Reflective Discrete Router}

The discrete method is trained and deployed independently; it is neither an
ensemble member nor a second stage of the soft-prompt router. A candidate $P$
combines an instruction with demonstrations sampled across outcomes. The
2B generative router emits $r_P(x,q)$, while a separate 4B model reflects on
execution traces during GEPA optimization. The adapter records correctness,
realized cost, failure type, and $S_{\mathrm{cost}}(x,r_P)$. Feedback
distinguishes critical under-escalation, unnecessary Mini use, Nano-only
success, and unrecoverable \bw cases. GEPA proposes policy mutations and
retains candidates on the validation frontier.

Discrete selection maximizes mean validation route score, then prefers fewer
critical misroutes, invalid outputs, and unnecessary Mini calls, followed by
higher answer accuracy. A scored prompt first emits a raw value in
$[0,100]$. The parser then computes
\[
s_P(x)=\operatorname{clip}(\tilde{s}_P(x),0,100)/100\in[0,1]
\]
and predicts Mini when $s_P(x)\ge\tau$. Reported thresholds use the
normalized scale.
Matched comparisons use the same validation accuracy floor for both routers. 
